% SI numbering: first-citation order (tools/si_numbering.py, round 6, 2026-10-05)
% discussion.tex : 고찰(국문판). 영문판 paper/manuscript/en/sections/discussion.tex(2026-10-05 2차 수정판)를
% 문단과 문장 단위로 옮겼다. 수치, 인용 키, 자리표시 문자열은 영문판과 같다.

\section{고찰}\label{sec:discussion}

결과는 새 영구동토 지역을 위한 라벨 수별 워크플로를 정한다(그림~6d). 대상 라벨이 없을 때 기준은 원천 계수(또는 연도 정합) Stefan 식이며, ML은 물리 유사라벨이나 저가중 잔차로만 들어간다. 라벨 3--10개에서는 계수를 재보정하고(오차가 낮아진 지역은 4곳 가운데 1곳) 저가중 잔차를 더한다. 라벨 40개부터는 대상 라벨 안 교차검증으로 방법을 고르며, 그 이득은 캐나다에서 생겼다. 이 풀에서 라벨 3--160개에 권하는 방법의 점 추정 오차는 원천 계수 모델보다 높았다 [DECISION: Fig 7d 경로의 y 기준과 풀, 그림 명세 D-13]. 라벨이 수천 개인 지역 안에서는 잔차 ML이 뒤따르며, 그 이득은 기후 격자 셀 사이에 있다(그림~6a--c). 배치는 사전 지정한 절차로만 들어간다. 라벨을 흩어 놓는 것이 캐나다에서는 오차를 낮추었고 레나델타에서는 낮추지 못했기 때문이다(그림~7).

이득은 Stefan 계수 $E$의 수준 보정과 기후 격자 규모의 편향 보정에서 왔다. $E$는 토양 열전도도, n-인자, 토양 수분과 얼음의 잠열을 한데 묶은 계수이며\cite{Riseborough2008}, 그래서 지역 Stefan 지도는 토지 피복별 계수를 썼다\cite{Nelson1997,ShiklomanovNelson2002}. 이 연구에서 $E$는 융해 지수의 편향과 라벨 정의도 흡수한다(보충 주석~1). 라벨 전량에서 원천 계수에 대한 잔차 ML의 4지역 평균 감소 가운데 85\%는 재보정 몫이었고, 이 평균은 러시아 서부가 지배하였다(14.0~cm 가운데 재보정 13.5~cm; 사후 산술). 계수가 원천 값에 가장 가까웠으나 라벨 블록과 채점 블록 사이에서 달랐던 캐나다에서는 라벨 10개 재보정이 오차를 키웠고 선정 규칙이 대개 잔차 가중 1.0을 골랐으며, 우리는 이를 라벨 위치의 이질성을 보정한 것으로 해석한다.

라벨 0개 결과는 랜덤 포레스트가 Stefan 식보다 일반화가 나빴던 알래스카 무작위 분할 비교와 방향이 같다\cite{Gautam2025}. 앞선 알래스카 지역 내 분석 [DECISION: 대회 보고서 인용 여부]은 56개 설정 가운데 최선인 물리 잔차에서 이 연구의 더 엄격한 시험보다 큰 Stefan 식 대비 이득을 보고하였다(보충 주석~2). kNNDM\cite{Linnenbrink2024}은 이 군집 표본(알래스카 13,606행이 343개 위치에 있음)에서 지도와 라벨 사이의 거리를 흉내 낼 뿐이다.

다른 설명의 가능성은 낮다. Stefan 유사라벨은 섞은 유사라벨과 상수 유사라벨보다 오차가 낮았고(그림~3b), 선형 융해 지수 위약과의 차이는 0.5~cm 미만이었다. 따라서 이득은 제곱근 형태보다 셀 수준의 융해 지수 정보를 반영한다. 재보정 이득은 러시아 서부와 동부의 단년 라벨에서도 유지되었고, 캐나다에서는 점 라벨일 때 1~km 위치 평균일 때보다 0.86~cm 작았다(중앙값, 라벨 10개). 이 결과가 워크플로의 라벨 단위를 정한다. 물리식과 CatBoost 결과는 계산 환경 사이에서 일치하였고(최대 차이 $3.6\times10^{-14}$~cm) 신경망은 최대 5.6~cm 달랐으므로, 워크플로는 결정적인 이 방법들만 쓴다.

워크플로는 측정 지점의 순위를 매기지 않고 기후 격자 셀 안의 변동을 예측하지 않는다. 1~km 지도(보충 그림~S1--S5)는 표시 해상도이고, 그 정보 해상도는 기후 입력(약 10~km)과 토양 입력(약 5~km)에 묶인다. 결론은 지표 투과 레이더(ground-penetrating radar) 라벨을 포함한 원천 자료에 한정되며, 훈련 범위보다 따뜻한 기후는 평가할 수 없었다. 원천 자료의 알래스카 비중(원천 셀의 78--94\%)은 라벨이 적을 때 가장 크게 작용한다. 원천 계수가 재보정 계수에서 라벨 10개의 가중을 갖기 때문이다. 확인적 근거는 재사용한 주 지역 4곳과 얕은 레짐의 새 지역 1곳에서 온다. 비맹검인 티베트 고원은 유일한 깊은 레짐 지역이라 따로 보고하므로, 결과를 이 지역들 너머로 일반화하지 않는다. 값은 시험 설계에 따라 달라지므로(그림~1e), 무작위 분할 연구와는 직접 비교할 수 없다. 격자 셀 안의 토양 수분, 유기층 두께, 적설 기간, 미지형은 공변량에 거의 없으므로, 알래스카에서 설명한 격자 안 비율(1\% 미만)은 물리적 한계가 아니라 이 입력의 한계를 나타낸다. 현장 토양 수분도 격자 안 잔차를 설명하지 못했다(진단; 보충 표~S2). 방법 선정, 배치, 라벨이 많은 지역의 시험, 편향 상관, 이득 분해는 앞선 결과를 본 뒤 설계하였으므로 내부 사전 등록한 대비보다 약한 근거다.

각 단계에는 조건이 있다. 라벨은 1~km 위치 평균이고, 재보정은 라벨 10개 대비에 근거하며, 대상 라벨 안 교차검증 선정은 라벨이 덮는 0.5$^{\circ}$ 블록 수에 좌우되고 라벨 10개의 레나델타에서 고정 레시피보다 오차를 키웠다. 워크플로 전체는 같은 지역을 다시 무작위로 나눈 분할에서만 시험하였고, 그 배치 단계는 무작위 추출로 바뀌었으며, 비열등은 풀 평균에서만 성립하고 러시아 동부, 알래스카, 레나델타에서는 성립하지 않았다. 다음에 필요한 자료는 개발에 쓰지 않은 독립 지역 \textcolor{blue}{[PENDING: XF, new public regions; registered data deadline 11 October 2026, 14:22 KST; no eligible new macro region found so far]}, 같은 블록에서 채점한 기존 ALT 지도\cite{Ran2022a,Wei2026}(이 연구에서는 등록 이탈로 비교; 보충 표~S2), 그리고 무엇보다 기후 격자 셀 안에서 변하는 공변량이다. 알래스카에서 재보정 Stefan 식 오차 제곱의 72\%가 격자 셀 안에 있었다. \textcolor{blue}{[PENDING: XE-a, XE-b, satellite inputs (stage 2); registered deadline 8 October 2026, 14:22 KST]}
